\section{Theoretical Foundations}
\label{sec:theory}

A core innovation of WuYa is its grounding in classical theories from the philosophy of science, statistics, and the sociology of science. This section establishes the theoretical basis by tracing each of the five evaluation dimensions to its philosophical origins, demonstrating that the criteria we take for granted are not arbitrary but have deep historical roots.

\subsection{Evidence Confidence: From Falsifiability to Bayesian Updating}

The fundamental question of what makes a scientific claim credible has been central to philosophy of science since the 20th century. Karl Popper's \textit{The Logic of Scientific Discovery} (1934) introduced the principle of \textbf{falsifiability} as the demarcation criterion between science and non-science. A scientific theory must make predictions that can, in principle, be refuted by empirical evidence \citep{popper1959logic}.

Popper's naive falsificationism was subsequently refined by Imre Lakatos' \textbf{research programmes} methodology (1970), which argued that the unit of evaluation is not a single hypothesis but an entire research programme. A programme is ``progressive'' if it predicts novel facts, and ``degenerating'' if it merely accommodates existing facts \citep{lakatos1970falsification}. This has direct implications for paper evaluation: reviewers should assess not just isolated findings but whether the work contributes to a progressive research programme.

The mathematical framework for quantifying evidence support was provided by Thomas Bayes' theorem (1763), which enables dynamic updating of belief degrees based on new evidence \citep{bayes1763essay}. Modern Bayesian inference provides the foundation for p-values, confidence intervals, and Bayesian factors used in statistical evaluation.

\begin{table}[t]
\centering
\caption{Theoretical Origins of Evaluation Dimensions}
\label{tab:theory-mapping}
\begin{tabular}{@{}llll@{}}
\toprule
\textbf{Dimension} & \textbf{Primary Theory} & \textbf{Foundational Work} & \textbf{Agent Component} \\
\midrule
Evidence Confidence & Falsifiability + Bayesianism & Popper (1934), Bayes (1763) & Evidence Sub-agent \\
Topic Innovation & Paradigm Shift + Knowledge Cycles & Kuhn (1962), Mokyr (2002) & Innovation Sub-agent \\
Method Validity & Causal Inference + Validity Theory & Fisher (1925), Pearl (2000) & Method Sub-agent \\
Application Potential & Linear Model + TRL & Bush (1945), NASA (1970s) & Application Sub-agent \\
Scientific Norms & CUDOS Principles & Merton (1942) & CUDOS Sub-agent \\
\bottomrule
\end{tabular}
\end{table}

Table~\ref{tab:theory-mapping} presents the mapping from evaluation dimensions to their theoretical foundations and corresponding \subagent{} components in WuYa.

\subsection{Topic Innovation: From Schumpeter to Kuhn to Knowledge Cycles}

Innovation is the engine of scientific progress, but what constitutes ``innovation'' varies across disciplines and eras. Joseph Schumpeter's theory of economic development (1912) first defined innovation as ``new combinations of production factors'' \citep{schumpeter1934theory}, providing the economic foundation for understanding innovation as recombination rather than mere accumulation.

Thomas Kuhn's \textit{The Structure of Scientific Revolutions} (1962) revolutionized the understanding of scientific change by introducing the concept of \textbf{paradigms} and \textbf{paradigm shifts} \citep{kuhn1970structure}. Scientific progress is not cumulative but proceeds through ``normal science $\rightarrow$ crisis $\rightarrow$ revolution'' cycles. Research representing paradigm shifts constitutes the highest level of innovation, while filling gaps within an existing paradigm represents a lower level.

Joel Mokyr's \textit{The Gifts of Athena} (2002) further developed this framework by distinguishing between propositional knowledge (Q-knowledge: ``knowing what'') and prescriptive knowledge (A-knowledge: ``knowing how'') \citep{mokyr2002gifts}. The Industrial Revolution succeeded because institutions enabled free flow between Q and A knowledge, creating a self-reinforcing dual cycle. This has implications for evaluating innovation: the highest-value research bridges theoretical understanding with practical applicability.

Philippe Aghion and Peter Howitt's \textbf{creative destruction} model (1992) formalized how continuous innovation drives long-term growth by replacing outdated technologies \citep{aghion1992model}. This informs the evaluation criterion: research should not only introduce new ideas but create opportunities for future innovation.

\subsection{Method Validity: From Fisher to Campbell to Pearl}

Methodological rigor ensures that research can reliably answer scientific questions. Ronald Fisher's \textit{Statistical Methods for Research Workers} (1925) established the gold standard of experimental design, including randomized controlled trials, ANOVA, and p-values \citep{fisher1925statistical}. These methods remain foundational for quantitative research evaluation.

Donald Campbell and Julian Stanley's \textbf{internal and external validity} framework (1963) systematically catalogued threats to research reliability \citep{campbell1963experimental}. Internal validity concerns whether causal inferences are warranted; external validity concerns generalizability. Eight threats to internal validity were identified: history, maturation, testing, instrumentation, statistical regression, selection, experimental mortality, and selection-maturation interaction.

Judea Pearl's \textbf{causal inference theory} (2000) provided a unified mathematical framework for causal reasoning, including causal graphs and do-calculus \citep{pearl2009causality}. The principle that ``no causation without a causal model'' has become central to modern methodological evaluation.

\subsection{Application Potential: From Bush to NASA TRL}

The practical value of research has been a concern of science policy since Vannevar Bush's \textit{Science: The Endless Frontier} (1945), which proposed the linear model of innovation: basic research $\rightarrow$ applied research $\rightarrow$ development $\rightarrow$ commercialization \citep{bush1945science}. While this linear model has been critiqued as oversimplified, it established the principle that basic research is the pacemaker of technological progress.

NASA's \textbf{Technology Readiness Level} (TRL) scale (1970s) provided a quantifiable framework for assessing technological maturity across nine levels, from basic principles observation (TRL 1) to actual system operational deployment (TRL 9) \citep{nasa1995trl}. This enables evaluation of whether research has clear pathways to practical application.

Mokyr's bidirectional knowledge interaction model enriches this picture by emphasizing that industrial translation is not one-way but involves feedback from application back to basic research \citep{mokyr2002gifts}.

\subsection{Scientific Norms: Merton's CUDOS Principles}

Robert Merton's normative structure of science (1942) established the ethical foundation for academic evaluation through the CUDOS principles \citep{merton1973normative}:

\begin{itemize}
    \item \textbf{Communism}: Scientific knowledge is a public good, to be shared openly.
    \item \textbf{Universalism}: Evaluation criteria are objective, independent of researcher identity.
    \item \textbf{Disinterestedness}: Research aims at truth, not personal gain.
    \item \textbf{Organized Skepticism}: All claims must undergo critical scrutiny by the scientific community.
\end{itemize}

These norms form the ethical bedrock of peer review, ensuring that evaluation is fair, transparent, and oriented toward collective knowledge advancement rather than individual or institutional interests.

\subsection{From Principles to Architecture: The Mapping Methodology}

WuYa's design philosophy is to map each theoretical principle to a corresponding architectural component:

\begin{enumerate}
    \item \textbf{Separation of concerns}: Different theoretical systems correspond to independent \subagents{}, avoiding prompt interference and enabling independent iteration.
    \item \textbf{Layered knowledge injection}: ``Principles'' (dao) are internalized in system prompts as evaluation frameworks; ``evidence'' (shu) is retrieved on-demand through \rag{} for concrete justification.
    \item \textbf{Gatekeeping priority}: CUDOS evaluation precedes expert evaluation, reflecting the priority of scientific norms over substantive assessment.
\end{enumerate}

This mapping ensures that WuYa's evaluations are not just technically correct but theoretically grounded---each judgment can be traced back to established principles from the philosophy of science.
